% Subsection 2.4.1: 可执行正确性

可执行正确性是代码生成最核心、最严格的评价标准，它要求模型输出代码能够在给定测试用例下\textbf{无异常编译、正常运行、输出完全符合预期结果}。

与 BLEU、ROUGE、CodeBLEU 等基于文本匹配的指标不同，可执行正确性不关注代码表面形式，而关注\textbf{运行时行为与逻辑正确性}，更能反映算法代码的实际可用性。其评价流程通常包含以下环节：
\begin{itemize}
    \item 语法校验：代码符合目标编程语言格式规范，无编译错误；
    \item 运行校验：代码可正常启动执行，不出现崩溃、死循环或超时；
    \item 结果校验：在标准测试用例上输入输出完全匹配标准答案；
    \item 边界校验：对特殊用例、极端条件能够正确处理，不出现逻辑漏洞。
\end{itemize}
在面向算法题的代码生成中，可执行正确性是判断方案有效性的首要依据，也是本文所有实验的基础评价准则。